\documentclass[a4paper,12pt]{article}
\usepackage[utf8]{inputenc}
\usepackage[T1]{fontenc}
\usepackage[ngerman]{babel}
\usepackage{amsmath}
\usepackage{amssymb}
\usepackage{enumitem}
\usepackage{geometry}
\geometry{a4paper, margin=2.5cm}

\title{Einsendeaufgaben -- Lektion 2\\[0.5em]
\large Modul 61111: Mathematische Grundlagen}
\author{}
\date{}

\begin{document}
\maketitle

\section*{Aufgabe 2.3}

\begin{enumerate}[label=\alph*)]

\item Sei $\lambda \in K \setminus \{0\}$ und $v, w \in V$ beliebig. Setze
\[
  x := \frac{v - w}{\lambda}.
\]
Dann folgt:
\[
  \lambda x + w = \lambda \cdot \frac{v - w}{\lambda} + w = v - w + w = v. \qquad \checkmark
\]

\item Gilt \textbf{nicht}. Gegenbeispiel:
\[
  V = \mathbb{R}^2, \quad v = \begin{pmatrix}1\\0\end{pmatrix}, \quad w = \begin{pmatrix}0\\1\end{pmatrix}.
\]
Es ist $\lambda v = \begin{pmatrix}\lambda\\0\end{pmatrix} \neq \begin{pmatrix}0\\1\end{pmatrix}$ für alle $\lambda \in \mathbb{R}$.

\item Sei $\lambda \in K \setminus \{0\}$ und $v \in V$ mit $\lambda v = 0$ beliebig.
Wegen $\lambda \neq 0$ folgt:
\[
  v = \frac{0}{\lambda} = 0.
\]

\end{enumerate}

\end{document}
